\section{Probabilidade Condicional e Independência}

Em experimentos aleatórios, um determinado evento pode ser mais ou menos frequente a
depender do conhecimento da ocorrência de outro evento. Nesse caso, dizemos que um evento
está condicionado a outro.

\begin{def:probabilidade condicional}
Seja $(\Omega, \mathcal{B}, P)$ um espaço de probabilidade.
Sejam $E$ e $F$ eventos de $\mathcal{B}$.
A probablidade condicional do evento $E$ dado $F$, denotada por $P(E|F)$, é dada por
\[
	P(E|F) = \frac{P(E \cap F)}{P(F)}
\]
\end{def:probabilidade condicional}

Quando a interseção entre dois eventos é vazia, então teremos que a probablidade
condicional será $0$.
A probablidade condicional induz uma forma de determinar a probabilidade do evento
dado como interseção de vários outros.

\begin{prop:regra da multiplicacao}[Regra do Produto]
Seja $\mathcal{F} = \{E_i\}_n$ uma família de eventos disjuntos de um
espaço de probabilidade $(\Omega, \mathcal{B}, P)$.
\[
	P\left(\bigcap_{E \in \mathcal{F}} E\right) =
	\prod_{i=1}^n P\left(E_i\ \middle|\ \bigcap_{j < i} E_j\right)
\]
\begin{proof}
	A prova será por indução em $n$

	Base: $n = 1$ é trivial.

	Hipótse de Indução: admita que o resultado vale para uma família de $n$ eventos. Se
	$|\mathcal{F}| = n + 1$, então façamos
	\[
		\bigcap_{j < n+1} E_j  = E'
	\]
	e obtemos que
	\[
		P\left(\bigcap_{E \in \mathcal{F}} E\right) = P(E' \cap E_n) = P(E_n | E') P(E').
	\]
	Como a hipótese de indução vale para $E'$, então
	\[
		P\left(\bigcap_{E \in \mathcal{F}} E\right) =
		P\left(E_n\ \middle|\ \bigcap_{j < n+1} E_j\right) \cdot
		\prod_{i=1}^n P\left(E_i\ \middle|\ \bigcap_{j < i} E_j\right) =
		\prod_{i=1}^{n+1} P\left(E_i\ \middle|\ \bigcap_{j < i} E_j\right)
	\]
	Uma vez que a proposição vale para $n+1$, então vale para todo $n \in \mathbb{N}$
	pelo princípio de indução.
\end{proof}
\end{prop:regra da multiplicacao}

Eventos independentes são aqueles cuja probabilidade de ocorrência não depende
da ocorrência de outro. Isso equivale a dizer que, o conhecimento do acontecimento
de um evento não muda a crença na ocorrência de outro.
\begin{def:independencia}
Dois eventos $E$ e $F$ de um espaço $(\Omega, \mathcal{B}, P)$ são independentes se, e
somente se
\[
	P(E \cap F) = P(E) \cdot P(F)
\]
\end{def:independencia}

\begin{prop:independencia}
Se $E$ e $F$ são eventos independentes de $(\Omega, \mathcal{B}, P)$, então
\[
	P(E | F) = P(E) \quad \mathit{ e }\quad P(F | E) = P(F)
\]
\begin{proof}
	Por definição de probabilidade condicional,
	\[
		P(E \cap F) = P(E|F) \cdot P(F)
	\]
	mas como $E$ e $F$ são independentes por hipótese,
	\begin{gather*}
		P(E) \cdot P(F) = P(E|F) \cdot P(F)\\
		P(E) = P(E|F)
	\end{gather*}
	e o mesmo pode ser mostrado analogamente para $P(F)$.
\end{proof}
\end{prop:independencia}

\begin{prop:independencia complementar}
\label{prop:independencia complementar}
Sejam $E$ e $F$ eventos em $(\Omega, \mathcal{B}, P)$.
Se eles são independentes, então os seguintes eventos também o são:
\begin{enumerate}
	\item $E$ e $F^c$
	\item $E^c$ e $F$
	\item $E^c$ e $F^c$
\end{enumerate}
\begin{proof}
	Para (1), é verdade que
	\begin{align*}
		P(E \cap F^c) & = P(E) - P(E \cap F) \\
		              & = P(E) [1 - P(F)]    \\
		              & = P(E) P(F^c)
	\end{align*}
	A expressão (2) é análoga a (1). Por fim, para expressão (3), conclua que
	\begin{align*}
		P(E^c \cap F^c) & = P((E \cup F)^c) = 1 - P(E \cup F) \\
		                & = 1 - P(E) - P(F) + P(E \cap F)     \\
		                & = P(E^c) - P(F) + P(E) \cdot P(F)   \\
		                & = P(E^c) - P(F) [1 - P(E)]          \\
		                & = P(E^c) - P(F) \cdot P(E^c)        \\
		                & = P(E^c) [1 - P(F)]                 \\
		                & = P(E^c) \cdot P(F^c)               \\
	\end{align*}
\end{proof}
\end{prop:independencia complementar}

\begin{thm:probabilidade total}[Teorema de Probablidade Total]
Seja o espaço de probabilidade $(\Omega, \mathcal{B}, P)$.
Seja $\mathcal{F} = \{F_i\}_\mathbb{N}$, com $\mathcal{F} \subset \mathcal{B}$, uma partição de
$\Omega$.
Se $E \in \mathcal{B}$, então
\[
	P(E) = \sum_{F \in \mathcal{F}} P(E) \cdot P(E | F)
\]
\begin{proof}
	Para quaisquer $i, j \in \mathbb{N}$ temos que
	\[
		(E \cap F_i) \cap (E \cap F_j)  = \varnothing.
	\]
	Com isso, podemos usar o Axioma III de Kolmogorov,
	\[
		P\left(\bigcup_{F \in \mathcal{F}} E \cap F_i \right) = \sum_{F \in \mathcal{F}} P(E \cap F).
	\]
	Em vista de que
	\[
		\bigcup_{F \in \mathcal{F}} E \cap F_i = E \cap \Omega = E,
	\]
	então
	\[
		P(E) = \sum_{F \in \mathcal{F}} P(E \cap F),
	\]
	e pela Proposição \ref{prop:independencia complementar}, conclui-se que
	\[
		P(E) = \sum_{F \in \mathcal{F}} P(E) \cdot P(E | F).
	\]
\end{proof}
\end{thm:probabilidade total}
